\documentclass[11pt]{article}
\usepackage{preamble}
\setlength{\parskip}{4pt}

\begin{document}

\daybanner{AP Statistics \textperiodcentered\ Unit 4 \textperiodcentered\ Topic 4.4 \textperiodcentered\ B: 2027-01-19 \textperiodcentered\ E: 2027-03-05 \textperiodcentered\ TEACHER KEY}{Two Measurements per Technician}

\sectionbanner{CED TETHER}

{\small
\begin{itemize}[leftmargin=1.5em,itemsep=2pt,topsep=2pt]
  \item \textbf{Skill 1.E} --- Identify an appropriate inference method for significance tests.
  \item \textbf{Skill 1.F} --- Identify null and alternative hypotheses.
  \item \textbf{Skill 4.C} --- Verify that inference procedures apply in a given situation.
  \item \textbf{EU VAR-7} --- The t-distribution may be used to model variation.
  \item \textbf{LO VAR-7.B} --- Identify an appropriate testing method for a population mean with unknown $\sigma$, including the mean difference between values in matched pairs. [Skill 1.E]
  \item \textbf{EK VAR-7.B.1} --- The appropriate test for a population mean with unknown $\sigma$ is a one-sample t-test for a population mean.
  \item \textbf{EK VAR-7.B.2} --- Matched pairs can be thought of as one sample of pairs. Once differences between pairs of values are found, inference for significance testing proceeds as for a population mean.
  \item \textbf{LO VAR-7.C} --- Identify the null and alternative hypotheses for a population mean with unknown $\sigma$, including the mean difference between values in matched pairs. [Skill 1.F]
  \item \textbf{EK VAR-7.C.1} --- The null hypothesis for a one-sample t-test for a population mean is $H_0$: $\mu$ = $\mu_0$, where $\mu_0$ is the hypothesized value. Depending upon the situation, the alternative hypothesis is $H_a$: $\mu$ < $\mu_0$, or $H_a$: $\mu$ > $\mu_0$, or $H_a$: $\mu$ $\neq$ $\mu_0$.
  \item \textbf{EK VAR-7.C.2} --- When finding the mean difference, $\mu_D$, between values in a matched pair, it is important to define the order of subtraction.
  \item \textbf{LO VAR-7.D} --- Verify the conditions for the test for a population mean, including the mean difference between values in matched pairs. [Skill 4.C]
  \item \textbf{EK VAR-7.D.1} --- In order to make statistical inferences when testing a population mean, we must check for independence and that the sampling distribution is approximately normal: a. To check for independence: (i) Data should be collected using a random sample or a randomized experiment. (ii) When sampling without replacement, check that n $\leq$ 10\% N.
  \item b. To check that the sampling distribution of $\bar{x}$ is approximately normal (shape): (i) If the observed distribution is skewed, n should be greater than 30. (ii) If the sample size is less than 30, the distribution of the sample data should be free from strong skewness and outliers.
\end{itemize}
}
\textbf{Essential question:} How does the study structure determine the parameter, hypotheses, and conditions for a mean test?\par
\textbf{DOK-3 skill:} 4.C \textperiodcentered\ \textbf{FRQ pattern:} {\small\texttt{paired-\allowbreak{}test-\allowbreak{}setup-\allowbreak{}and-\allowbreak{}consequences}}\par
\textbf{DOK rationale:} Part (c) coordinates repeated measures, subtraction order, tail direction, independent sample size, and shape evidence, then traces the uncertainty and research-question consequences of the rejected setup.\par

\frameworkphaseheader{Do Now (first take) $\rightarrow$ Explore (video + follow-along) $\rightarrow$ Exit (finish + turn in)}{1 $\rightarrow$ 3}{45}{%
\begin{itemize}[leftmargin=1.5em,itemsep=2pt,topsep=2pt]
  \item Display the board and ask for one sentence using the study description.
  \item Begin the 7.4 video follow-along at minute 5.
  \item Return to the board at minute 33. Students finish the shortened parts and justify their judgment.
  \item Collect at the door. Score part (c) only using the three required elements.
\end{itemize}
}{%
\begin{itemize}[leftmargin=1.5em,itemsep=2pt,topsep=2pt]
  \item Write a one-sentence first take before the video.
  \item Complete the video follow-along worksheet.
  \item Finish (a)--(c), revise the first take in the exit line, and turn in the sheet.
\end{itemize}
}{%
\begin{itemize}[leftmargin=1.5em,itemsep=2pt,topsep=2pt]
  \item Which specific number or study detail supports your choice?
  \item What claim would go beyond the evidence, and why?
\end{itemize}
}{Circulate and ask for evidence. Accept a concise response; do not require omitted calculations or a second full inference write-up.}

\begin{callout}{\IconWarn\ WATCH FOR}{calloutred}
\begin{itemize}[leftmargin=1.5em,itemsep=2pt,topsep=2pt]
  \item Choosing a speaker without a reason.
  \item Copying numbers without explaining what they support.
\end{itemize}
\end{callout}

\sectionbanner{SCORING PART (c) --- E / P / I}

\textbf{Expected elements}
\begin{itemize}[leftmargin=1.5em,itemsep=2pt,topsep=2pt]
  \item \textbf{paired-units} (required) --- Justifies a one-sample test of the 18 within-technician differences
  \item \textbf{correct-hypotheses} (required) --- Uses the population mean difference, zero null, and lower-tail alternative
  \item \textbf{rejects-independent-times} (required) --- Explains that the two times from each person are linked
\end{itemize}

{\small\renewcommand{\arraystretch}{1.25}
\noindent\begin{tabular}{|p{0.5in}|p{5.9in}|}
\hline
\textbf{E} & Justifies the paired units and explains the lower-tail population hypotheses using new minus old and the reduction question. \\ \hline
\textbf{P} & Shows substantive valid reasoning for at least one required element, but does not meet all E requirements. Credit correct reasoning even when another claim is wrong. \\ \hline
\textbf{I} & Shows no substantive valid reasoning for any required element; a name, unsupported choice, or copied number alone does not count. \\ \hline
\end{tabular}}

\clearpage
\focusbanner{TODAY'S PROBLEM --- annotated}

Suppose a company takes an SRS of 18 of its 300 technicians. Each completes the same task with old and new interfaces; interface order is randomly assigned. The plot shows $d=\text{new time}-\text{old time}$. Researchers ask whether the new interface reduces the population mean time.\par\smallskip \textbf{Nia:} ``The technician seems to be the independent unit, so I would build one difference per person; the subtraction order matters for the tail.''\par \textbf{Colin:} ``There are 36 recorded times, so I would use two independent samples of 18 and a two-sided alternative.''\par
\begin{center}
\begin{tikzpicture}
\begin{axis}[
  width=0.61\linewidth,
  height=1.15in,
  ymin=0, ymax=2, ytick=\empty, axis y line=none, axis x line=bottom,
  xlabel={New-interface time minus old-interface time (seconds)}, tick label style={font=\small}, label style={font=\small}
]
\addplot[only marks, mark=*, mark size=1.7pt, color=dayblue] coordinates {
  (-4.2,1)
  (-3.6,1)
  (-3.2,1)
  (-2.9,1)
  (-2.6,1)
  (-2.4,1)
  (-2.2,1)
  (-2.0,1)
  (-1.9,1)
  (-1.7,1)
  (-1.6,1)
  (-1.4,1)
  (-1.2,1)
  (-1.0,1)
  (-0.7,1)
  (-0.4,1)
  (0.0,1)
  (0.6,1)
};
\end{axis}
\end{tikzpicture}
\end{center}
\begin{firsttakebox}{FIRST TAKE (5 min)}
Would you compare each technician's two times or treat them as separate people? Explain in one sentence.\par
\answer{Ungraded. Everyday language is enough before the video; formal vocabulary belongs in the finish.}
\end{firsttakebox}

\textit{Rules callout ``SET UP A TEST FOR A MEAN (keep on the page)'' is printed on the student sheet.}\par\medskip

\dokbadge{1}~\textbf{(a)} Define $\mu_d$ in context.\par
\answer{$\mu_d$ is the mean new-minus-old task time for all 300 technicians.}

\dokbadge{2}~\textbf{(b)} Check the random-sample, 10 percent, and small-sample shape conditions in three short notes.\par
\answer{The 18 technicians were randomly sampled. $18\leq0.10(300)=30$. Since $n=18<30$, inspect the differences: the plot shows no strong skewness or outliers.}

\dokbadge{3}~\textbf{(c)} In 3--4 sentences, choose a setup and justify it. State the hypotheses using the given subtraction order and reduction question, then explain why the 36 times are not 36 independent people.\par
\answer{Nia is justified: use a matched-pairs one-sample $t$-test on the 18 differences. For new minus old, a reduction is negative, so $H_0:\mu_d=0$ and $H_a:\mu_d<0$. Each technician supplies two linked times, so there are 18 independent people, not 36. Treating the times as independent loses the pairing, and a two-sided alternative asks about any change rather than the planned reduction.}

\begin{summaryexitbox}{EXIT}
One thing the video changed about my first take: \hrulefill
\end{summaryexitbox}

\end{document}
